% GENERATED FILE. Edit canonical lessons, then run npm run generate:books.
% canonical-source-hash: fnv1a64:a06168a8bc9eb450
% canonical-lessons: JA-C137-ippai, JA-W137-pa, JA-C137-shinpai, JA-C137-enpitsu, JA-W137-pi, JA-C137-ippiki, JA-C137-kippu, JA-W137-pu, JA-C137-tenpura, JA-C137-perapera, JA-W137-pe, JA-R137-pencil-and-ticket

\chapter{Pencil and Ticket: Pa, Pi, Pu and Pe}
\label{ch:ja-pencil-and-ticket-pa-pi-pu-pe}
\hlchaptermodality{137}

\begin{chapteropening}
\textbf{By the end of this chapter:} \emph{I can write pa, pi, pu and pe by adding the small circle to signs I already write, and read words for things I want or need that use them, such as kippu, enpitsu and tenpura.}

Read ippai, shinpai, enpitsu, ippiki, kippu, tenpura and perapera, write the four new signs beside the signs they come from, write enpitsu and kippu from memory and say which one you show on a train, and write perapera.
\end{chapteropening}

\section[\texorpdfstring{\ja{いっぱい} (ippai)}{いっぱい (ippai)}]{\ja{いっぱい} (ippai) — full; a lot}

\label{lesson:JA-C137-ippai}



\begin{quote}
Three recalls before the new word.

\begin{itemize}
\raggedright
  \item \emph{Write it:} \textbf{\ja{ず}} from memory — \textbf{R2}, five lessons back
  \item \emph{Write it:} \textbf{\ja{び}} from memory — \textbf{R3}, twenty lessons back
  \item \emph{Recall:} say \emph{clever} — \textbf{R4}, eighty lessons back
\end{itemize}
\end{quote}

\subsection*{You'll want to know: \ja{いっぱい}}
\textbf{\ja{いっぱい}} — \emph{ippai} — \textquotedblleft{}full\textquotedblright{}, as a cup is full, or as you are after a big meal. It also means \textquotedblleft{}a lot\textquotedblright{}: a lot of people, a lot of work.

Four beats: \emph{i–p–pa–i}. \textbf{\ja{い}}, the small \textbf{\ja{っ}} that holds a beat, and the last \textbf{\ja{い}} you already write. The third sign is \textbf{\ja{は}} with the small circle at its upper right, and that circle makes a \emph{p}: \emph{ha} becomes \emph{pa}. The next lesson writes it.

\subsection*{Your turn}
\begin{itemize}
\raggedright
  \item \emph{Say it:} \emph{ippai}
  \item \emph{Say it:} \emph{ippai}, then count its beats aloud — four
  \item \emph{Recall:} say which sign in \textbf{\ja{いっぱい}} carries the small circle, and name the sign under it
\end{itemize}

\begin{tcolorbox}[breakable,colback=teal!4,colframe=teal!35!black,title={Before you move on}]
What is \textquotedblleft{}full\textquotedblright{}? (\textbf{\ja{いっぱい}}, \emph{ippai}.) How many beats? (\textbf{Four}.)
\end{tcolorbox}
\section[\texorpdfstring{\ja{ぱ} (pa)}{ぱ (pa)}]{\ja{ぱ} — \ja{は}, with the circle}

\label{lesson:JA-W137-pa}



\begin{quote}
Four recalls, then the sign.

\begin{itemize}
\raggedright
  \item \emph{Recall:} say \emph{full}, or \emph{a lot} — \textbf{R1}, one lesson back
  \item \emph{Recall:} say \emph{cool}, of a summer evening — \textbf{R2}, five lessons back
  \item \emph{Recall:} say \emph{Saturday} — \textbf{R3}, twenty lessons back
  \item \emph{Recall:} say \emph{foolish} — \textbf{R4}, eighty lessons back
\end{itemize}
\end{quote}

\subsection*{Script — the change}
\hlblockfigure{\detokenize{figures/JA-W137-pa-filmstrip.pdf}}{How \ja{ぱ} is written, stroke by stroke}

\begin{quote}
\textbf{\ja{は}} \emph{ha} + \textbf{\ja{゜}} $\to$ \textbf{\ja{ぱ}} \emph{pa}
\end{quote}

Write \textbf{\ja{は}} in full, in the order you learned it. Then add the small circle of the handakuten at the upper right, where the voicing mark would go. Keep it small and open, and do not let it touch the sign.

The circle works only on the \emph{h} row, and it always makes a \emph{p}: \emph{ha} becomes \emph{pa}. You saw it on \textbf{\ja{ほ}} for \textbf{\ja{ぽ}}, in \textbf{\ja{いっぽん}}.

With \textbf{\ja{いっ}} in front and \textbf{\ja{い}} after it, you can now write the word from the last lesson: \textbf{\ja{いっぱい}}, \emph{ippai}, full.

\subsection*{Your turn}
\begin{enumerate}
\raggedright
  \item Write \textbf{\ja{は}}, then \textbf{\ja{ば}}, then \textbf{\ja{ぱ}}, and say \emph{ha}, \emph{ba}, \emph{pa}.
  \item Write \textbf{\ja{ほ} \ja{ぼ} \ja{ぽ}} and \textbf{\ja{は} \ja{ば} \ja{ぱ}} in rows, and say each one.
  \item Write \textbf{\ja{いっぱい}} from memory and say \emph{ippai}.
\end{enumerate}

\begin{tcolorbox}[breakable,colback=teal!4,colframe=teal!35!black,title={Before you move on}]
Stop after one accurate form.

Source: \href{https://www.unicode.org/charts/PDF/U3040.pdf}{Unicode Hiragana chart}.
\end{tcolorbox}
\section[\texorpdfstring{\ja{しんぱい} (shinpai)}{しんぱい (shinpai)}]{\ja{しんぱい} (shinpai) — worry; worried}

\label{lesson:JA-C137-shinpai}



\begin{quote}
Four recalls before the new word.

\begin{itemize}
\raggedright
  \item \emph{Write it:} \textbf{\ja{ぱ}} from memory — \textbf{R1}, one lesson back
  \item \emph{Recall:} say \emph{a wind}, or \emph{a cold} — \textbf{R2}, five lessons back
  \item \emph{Recall:} say \emph{an entrance} — \textbf{R3}, twenty lessons back
  \item \emph{Recall:} say \emph{brave} — \textbf{R4}, eighty lessons back
\end{itemize}
\end{quote}

\subsection*{You'll want to know: \ja{しんぱい}}
\textbf{\ja{しんぱい}} — \emph{shinpai} — \textquotedblleft{}worry\textquotedblright{}: the feeling, or being worried about someone. It is the word in \emph{Don't worry}.

The new sign is spent at once. \textbf{\ja{し}} and \textbf{\ja{ん}}, then \textbf{\ja{ぱ}} and \textbf{\ja{い}}: four beats, \emph{shi–n–pa–i}. After \textbf{\ja{ん}} the \emph{p} comes out clearly, the lips closed first. You write every sign in it now.

\subsection*{Your turn}
\begin{itemize}
\raggedright
  \item \emph{Say it:} \emph{shinpai}
  \item \emph{Say it:} \emph{ippai}, then \emph{shinpai}, and listen for the \emph{pai} in both
  \item \emph{Write it:} \textbf{\ja{しんぱい}} from memory, with the circle small and open
\end{itemize}

\begin{tcolorbox}[breakable,colback=teal!4,colframe=teal!35!black,title={Before you move on}]
What does \ja{しんぱい} mean? (\textquotedblleft{}Worry; worried\textquotedblright{}.) Say it once more.
\end{tcolorbox}
\section[\texorpdfstring{\ja{えんぴつ} (enpitsu)}{えんぴつ (enpitsu)}]{\ja{えんぴつ} (enpitsu) — a pencil}

\label{lesson:JA-C137-enpitsu}



\begin{quote}
Four recalls before the new word.

\begin{itemize}
\raggedright
  \item \emph{Recall:} say \emph{worry} — \textbf{R1}, one lesson back
  \item \emph{Write it:} \textbf{\ja{ぜ}} from memory — \textbf{R2}, five lessons back
  \item \emph{Write it:} \textbf{\ja{ぐ}} from memory — \textbf{R3}, twenty lessons back
  \item \emph{Recall:} say \emph{fine}, small and detailed — \textbf{R4}, eighty lessons back
\end{itemize}
\end{quote}

\subsection*{You'll want to know: \ja{えんぴつ}}
\textbf{\ja{えんぴつ}} — \emph{enpitsu} — \textquotedblleft{}a pencil\textquotedblright{}, the one you may be writing with now.

Four beats: \emph{e–n–pi–tsu}. \textbf{\ja{え}}, \textbf{\ja{ん}} and \textbf{\ja{つ}} you already write. The third sign is \textbf{\ja{ひ}} with the small circle, so \emph{hi} becomes \emph{pi}. The next lesson writes it.

\subsection*{Your turn}
\begin{itemize}
\raggedright
  \item \emph{Say it:} \emph{enpitsu}
  \item \emph{Say it:} \emph{enpitsu}, then count its beats aloud — four
  \item \emph{Recall:} say which sign in \textbf{\ja{えんぴつ}} carries the small circle, and name the sign under it
\end{itemize}

\begin{tcolorbox}[breakable,colback=teal!4,colframe=teal!35!black,title={Before you move on}]
What is \textquotedblleft{}a pencil\textquotedblright{}? (\textbf{\ja{えんぴつ}}, \emph{enpitsu}.) How many beats? (\textbf{Four}.)
\end{tcolorbox}
\section[\texorpdfstring{\ja{ぴ} (pi)}{ぴ (pi)}]{\ja{ぴ} — \ja{ひ}, with the circle}

\label{lesson:JA-W137-pi}



\begin{quote}
Three recalls, then the sign.

\begin{itemize}
\raggedright
  \item \emph{Recall:} say \emph{a pencil} — \textbf{R1}, one lesson back
  \item \emph{Recall:} say \emph{an exit} — \textbf{R3}, twenty lessons back
  \item \emph{Recall:} say \emph{ugly} — \textbf{R4}, eighty lessons back
\end{itemize}
\end{quote}

\subsection*{Script — the change}
\hlblockfigure{\detokenize{figures/JA-W137-pi-filmstrip.pdf}}{How \ja{ぴ} is written, stroke by stroke}

\begin{quote}
\textbf{\ja{ひ}} \emph{hi} + \textbf{\ja{゜}} $\to$ \textbf{\ja{ぴ}} \emph{pi}
\end{quote}

Write \textbf{\ja{ひ}} in full: one stroke, a short rise, the deep bowl, and up again. Then add the small circle at the upper right, beside the end of the stroke. Keep it small and open.

The circle makes a \emph{p} again: \emph{hi} becomes \emph{pi}. The voicing mark on the same sign gives \textbf{\ja{び}}, \emph{bi}, so one base sign now has three readings.

With \textbf{\ja{えん}} in front and \textbf{\ja{つ}} after it, you can now write the word from the last lesson: \textbf{\ja{えんぴつ}}, \emph{enpitsu}, a pencil.

\subsection*{Your turn}
\begin{enumerate}
\raggedright
  \item Write \textbf{\ja{ひ}}, then \textbf{\ja{び}}, then \textbf{\ja{ぴ}}, and say \emph{hi}, \emph{bi}, \emph{pi}.
  \item Write \textbf{\ja{は} \ja{ぱ}} and \textbf{\ja{ひ} \ja{ぴ}} in pairs, and say each one.
  \item Write \textbf{\ja{えんぴつ}} from memory and say \emph{enpitsu}.
\end{enumerate}

\begin{tcolorbox}[breakable,colback=teal!4,colframe=teal!35!black,title={Before you move on}]
Stop after one accurate form.

Source: \href{https://www.unicode.org/charts/PDF/U3040.pdf}{Unicode Hiragana chart}.
\end{tcolorbox}
\section[\texorpdfstring{\ja{いっぴき} (ippiki)}{いっぴき (ippiki)}]{\ja{いっぴき} (ippiki) — one small animal}

\label{lesson:JA-C137-ippiki}



\begin{quote}
Four recalls before the new word.

\begin{itemize}
\raggedright
  \item \emph{Write it:} \textbf{\ja{ぴ}} from memory — \textbf{R1}, one lesson back
  \item \emph{Recall:} say \emph{full}, or \emph{a lot} — \textbf{R2}, five lessons back
  \item \emph{Recall:} say \emph{Monday} — \textbf{R3}, twenty lessons back
  \item \emph{Recall:} say \emph{cute} — \textbf{R4}, eighty lessons back
\end{itemize}
\end{quote}

\subsection*{You'll want to know: \ja{いっぴき}}
\textbf{\ja{いっぴき}} — \emph{ippiki} — \textquotedblleft{}one small animal\textquotedblright{}: one cat, one dog, one fish. Japanese counts animals with a word of their own, the way it counts long thin things with \textbf{\ja{いっぽん}}.

The new sign is spent at once. \textbf{\ja{い}}, a small \textbf{\ja{っ}}, \textbf{\ja{ぴ}} and \textbf{\ja{き}}: four beats, \emph{i–p–pi–ki}. The small \textbf{\ja{っ}} holds a beat before the \emph{p}, as it does in \textbf{\ja{いっぱい}}. You write every sign in it now.

\subsection*{Your turn}
\begin{itemize}
\raggedright
  \item \emph{Say it:} \emph{ippiki}
  \item \emph{Say it:} \emph{ippai}, then \emph{ippiki}, and listen for the held beat in both
  \item \emph{Write it:} \textbf{\ja{いっぴき}} from memory, with the \textbf{\ja{っ}} small and low
\end{itemize}

\begin{tcolorbox}[breakable,colback=teal!4,colframe=teal!35!black,title={Before you move on}]
What does \ja{いっぴき} mean? (\textquotedblleft{}One small animal\textquotedblright{}.) Say it once more.
\end{tcolorbox}
\section[\texorpdfstring{\ja{きっぷ} (kippu)}{きっぷ (kippu)}]{\ja{きっぷ} (kippu) — a ticket}

\label{lesson:JA-C137-kippu}



\begin{quote}
Four recalls before the new word.

\begin{itemize}
\raggedright
  \item \emph{Recall:} say \emph{one small animal}, one cat — \textbf{R1}, one lesson back
  \item \emph{Write it:} \textbf{\ja{ぱ}} from memory — \textbf{R2}, five lessons back
  \item \emph{Write it:} \textbf{\ja{げ}} from memory — \textbf{R3}, twenty lessons back
  \item \emph{Recall:} say \emph{persistent} — \textbf{R4}, eighty lessons back
\end{itemize}
\end{quote}

\subsection*{You'll want to know: \ja{きっぷ}}
\textbf{\ja{きっぷ}} — \emph{kippu} — \textquotedblleft{}a ticket\textquotedblright{}, above all the one you buy for a train.

Three beats: \emph{ki–p–pu}. \textbf{\ja{き}} and the small \textbf{\ja{っ}} you already write. The last sign is \textbf{\ja{ふ}} with the small circle, so \emph{fu} becomes \emph{pu}. The next lesson writes it.

\subsection*{Your turn}
\begin{itemize}
\raggedright
  \item \emph{Say it:} \emph{kippu}
  \item \emph{Say it:} \emph{kippu}, then count its beats aloud — three
  \item \emph{Recall:} say which sign in \textbf{\ja{きっぷ}} carries the small circle, and name the sign under it
\end{itemize}

\begin{tcolorbox}[breakable,colback=teal!4,colframe=teal!35!black,title={Before you move on}]
What is \textquotedblleft{}a ticket\textquotedblright{}? (\textbf{\ja{きっぷ}}, \emph{kippu}.) How many beats? (\textbf{Three}.)
\end{tcolorbox}
\section[\texorpdfstring{\ja{ぷ} (pu)}{ぷ (pu)}]{\ja{ぷ} — \ja{ふ}, with the circle}

\label{lesson:JA-W137-pu}



\begin{quote}
Four recalls, then the sign.

\begin{itemize}
\raggedright
  \item \emph{Recall:} say \emph{a ticket} — \textbf{R1}, one lesson back
  \item \emph{Recall:} say \emph{worry} — \textbf{R2}, five lessons back
  \item \emph{Recall:} say \emph{well}, as in healthy — \textbf{R3}, twenty lessons back
  \item \emph{Recall:} say \emph{noisy} — \textbf{R4}, eighty lessons back
\end{itemize}
\end{quote}

\subsection*{Script — the change}
\hlblockfigure{\detokenize{figures/JA-W137-pu-filmstrip.pdf}}{How \ja{ぷ} is written, stroke by stroke}

\begin{quote}
\textbf{\ja{ふ}} \emph{fu} + \textbf{\ja{゜}} $\to$ \textbf{\ja{ぷ}} \emph{pu}
\end{quote}

Write \textbf{\ja{ふ}} in full, in the order you learned it. Then add the small circle at the upper right, above the dot on the right. Keep it small and open.

The sign is read \emph{pu}. The \emph{f} of \emph{fu} goes, and the lips close for a \emph{p}. With the voicing mark the same sign gives \textbf{\ja{ぶ}}, \emph{bu}.

With \textbf{\ja{きっ}} in front of it, you can now write the word from the last lesson: \textbf{\ja{きっぷ}}, \emph{kippu}, a ticket.

\subsection*{Your turn}
\begin{enumerate}
\raggedright
  \item Write \textbf{\ja{ふ}}, then \textbf{\ja{ぶ}}, then \textbf{\ja{ぷ}}, and say \emph{fu}, \emph{bu}, \emph{pu}.
  \item Write \textbf{\ja{ぱ}}, \textbf{\ja{ぴ}}, \textbf{\ja{ぷ}} and \textbf{\ja{ぽ}}, and say each one.
  \item Write \textbf{\ja{きっぷ}} from memory and say \emph{kippu}.
\end{enumerate}

\begin{tcolorbox}[breakable,colback=teal!4,colframe=teal!35!black,title={Before you move on}]
Stop after one accurate form.

Source: \href{https://www.unicode.org/charts/PDF/U3040.pdf}{Unicode Hiragana chart}.
\end{tcolorbox}
\section[\texorpdfstring{\ja{てんぷら} (tenpura)}{てんぷら (tenpura)}]{\ja{てんぷら} (tenpura) — tempura, food fried in a light batter}

\label{lesson:JA-C137-tenpura}



\begin{quote}
Four recalls before the new word.

\begin{itemize}
\raggedright
  \item \emph{Write it:} \textbf{\ja{ぷ}} from memory — \textbf{R1}, one lesson back
  \item \emph{Recall:} say \emph{a pencil} — \textbf{R2}, five lessons back
  \item \emph{Recall:} say \emph{a bank} — \textbf{R3}, twenty lessons back
  \item \emph{Recall:} say \emph{shameless} — \textbf{R4}, eighty lessons back
\end{itemize}
\end{quote}

\subsection*{You'll want to know: \ja{てんぷら}}
\textbf{\ja{てんぷら}} — \emph{tenpura} — \textquotedblleft{}tempura\textquotedblright{}: fish, prawns or vegetables, dipped in a light batter and fried. English borrowed the word from Japanese.

The new sign is spent at once. \textbf{\ja{て}} and \textbf{\ja{ん}}, then \textbf{\ja{ぷ}} and \textbf{\ja{ら}}: four beats, \emph{te–n–pu–ra}. After \textbf{\ja{ん}} the \emph{p} comes out clearly, as in \textbf{\ja{しんぱい}}. You write every sign in it now.

\subsection*{Your turn}
\begin{itemize}
\raggedright
  \item \emph{Say it:} \emph{tenpura}
  \item \emph{Say it:} \emph{kippu}, then \emph{tenpura}, and listen for the \emph{pu} in both
  \item \emph{Write it:} \textbf{\ja{てんぷら}} from memory, four signs for four beats
\end{itemize}

\begin{tcolorbox}[breakable,colback=teal!4,colframe=teal!35!black,title={Before you move on}]
What does \ja{てんぷら} mean? (\textquotedblleft{}Tempura\textquotedblright{}, fried in batter.) Say it once more.
\end{tcolorbox}
\section[\texorpdfstring{\ja{ぺらぺら} (perapera)}{ぺらぺら (perapera)}]{\ja{ぺらぺら} (perapera) — fluent, fluently}

\label{lesson:JA-C137-perapera}



\begin{quote}
Four recalls before the new word.

\begin{itemize}
\raggedright
  \item \emph{Recall:} say \emph{tempura} — \textbf{R1}, one lesson back
  \item \emph{Write it:} \textbf{\ja{ぴ}} from memory — \textbf{R2}, five lessons back
  \item \emph{Write it:} \textbf{\ja{ぎ}} from memory — \textbf{R3}, twenty lessons back
  \item \emph{Recall:} say \emph{reliable} — \textbf{R4}, eighty lessons back
\end{itemize}
\end{quote}

\subsection*{You'll want to know: \ja{ぺらぺら}}
\textbf{\ja{ぺらぺら}} — \emph{perapera} — \textquotedblleft{}fluent, fluently\textquotedblright{}: it says that someone speaks a language with ease. It is one of the many doubled words Japanese uses for a sound or a manner: \emph{pera}, and \emph{pera} again.

Four beats: \emph{pe–ra–pe–ra}. \textbf{\ja{ら}} you already write. The other sign, twice, is \textbf{\ja{へ}} with the small circle, so \emph{he} becomes \emph{pe}. The next lesson writes it.

\subsection*{Your turn}
\begin{itemize}
\raggedright
  \item \emph{Say it:} \emph{perapera}
  \item \emph{Say it:} \emph{perapera}, then count its beats aloud — four
  \item \emph{Recall:} say which sign in \textbf{\ja{ぺらぺら}} carries the small circle, and name the sign under it
\end{itemize}

\begin{tcolorbox}[breakable,colback=teal!4,colframe=teal!35!black,title={Before you move on}]
What is \textquotedblleft{}fluent\textquotedblright{}? (\textbf{\ja{ぺらぺら}}, \emph{perapera}.) How many beats? (\textbf{Four}.)
\end{tcolorbox}
\section[\texorpdfstring{\ja{ぺ} (pe)}{ぺ (pe)}]{\ja{ぺ} — \ja{へ}, with the circle}

\label{lesson:JA-W137-pe}



\begin{quote}
Three recalls, then the sign.

\begin{itemize}
\raggedright
  \item \emph{Recall:} say \emph{fluent}, of someone who speaks with ease — \textbf{R1}, one lesson back
  \item \emph{Recall:} say \emph{one small animal}, one cat — \textbf{R2}, five lessons back
  \item \emph{Recall:} say \emph{pure white} — \textbf{R4}, eighty lessons back
\end{itemize}
\end{quote}

\subsection*{Script — the change}
\hlblockfigure{\detokenize{figures/JA-W137-pe-filmstrip.pdf}}{How \ja{ぺ} is written, stroke by stroke}

\begin{quote}
\textbf{\ja{へ}} \emph{he} + \textbf{\ja{゜}} $\to$ \textbf{\ja{ぺ}} \emph{pe}
\end{quote}

Write \textbf{\ja{へ}} in full: one stroke, up to the low peak and down to the right. Then add the small circle at the upper right, above the long way down. Keep it small and open.

The circle makes a \emph{p}: \emph{he} becomes \emph{pe}. With it you write every sign of the \emph{p} row: \textbf{\ja{ぱ} \ja{ぴ} \ja{ぷ} \ja{ぺ} \ja{ぽ}}.

With \textbf{\ja{ら}} after it, twice over, you can now write the word from the last lesson: \textbf{\ja{ぺらぺら}}, \emph{perapera}, fluent.

\subsection*{Your turn}
\begin{enumerate}
\raggedright
  \item Write \textbf{\ja{へ}}, then \textbf{\ja{べ}}, then \textbf{\ja{ぺ}}, and say \emph{he}, \emph{be}, \emph{pe}.
  \item Write \textbf{\ja{ぱ} \ja{ぴ} \ja{ぷ} \ja{ぺ} \ja{ぽ}} in order, and say each one.
  \item Write \textbf{\ja{ぺらぺら}} from memory and say \emph{perapera}.
\end{enumerate}

\begin{tcolorbox}[breakable,colback=teal!4,colframe=teal!35!black,title={Before you move on}]
Stop after one accurate form.

Source: \href{https://www.unicode.org/charts/PDF/U3040.pdf}{Unicode Hiragana chart}.
\end{tcolorbox}
\section[(dialogue)]{Pencil and ticket — the p row, complete}

\label{lesson:JA-R137-pencil-and-ticket}



\begin{quote}
Four recalls, then the review.

\begin{itemize}
\raggedright
  \item \emph{Write it:} \textbf{\ja{ぺ}} from memory — \textbf{R1}, one lesson back
  \item \emph{Recall:} say \emph{a ticket} — \textbf{R2}, five lessons back
  \item \emph{Recall:} say \emph{a family} — \textbf{R3}, twenty lessons back
  \item \emph{Recall:} say \emph{pitch black} — \textbf{R4}, eighty lessons back
\end{itemize}
\end{quote}

\subsection*{Your turn — read the seven words}
\emph{Read it:} each one below aloud, then say what it means

Read these aloud:

\begin{itemize}
\raggedright
  \item \textbf{\ja{いっぱい}} — full; a lot
  \item \textbf{\ja{しんぱい}} — worry
  \item \textbf{\ja{えんぴつ}} — a pencil
  \item \textbf{\ja{いっぴき}} — one small animal
  \item \textbf{\ja{きっぷ}} — a ticket
  \item \textbf{\ja{てんぷら}} — tempura
  \item \textbf{\ja{ぺらぺら}} — fluent
\end{itemize}

Each of these words holds one of the four signs this chapter wrote. With them, and \textbf{\ja{ぽ}} from chapter 18, the book writes every sign of the \emph{p} row, \textbf{\ja{ぱ} \ja{ぴ} \ja{ぷ} \ja{ぺ} \ja{ぽ}}.

\subsection*{Your turn — write}
\begin{enumerate}
\raggedright
  \item \emph{Write it:} \textbf{\ja{は} \ja{ぱ}}, \textbf{\ja{ひ} \ja{ぴ}}, \textbf{\ja{ふ} \ja{ぷ}} and \textbf{\ja{へ} \ja{ぺ}}, in pairs
  \item \emph{Write it:} \textbf{\ja{えんぴつ}} and \textbf{\ja{きっぷ}} from memory, and say which one you show on a train
  \item Say \emph{perapera}. \emph{Write it:} \textbf{\ja{ぺらぺら}}
\end{enumerate}

\begin{tcolorbox}[breakable,colback=teal!4,colframe=teal!35!black,title={Before you move on}]
Full? (\textbf{\ja{いっぱい}}, \emph{ippai}.) Worry? (\textbf{\ja{しんぱい}}, \emph{shinpai}.) A pencil? (\textbf{\ja{えんぴつ}}, \emph{enpitsu}.) One small animal? (\textbf{\ja{いっぴき}}, \emph{ippiki}.) A ticket? (\textbf{\ja{きっぷ}}, \emph{kippu}.) Tempura? (\textbf{\ja{てんぷら}}, \emph{tenpura}.) Fluent? (\textbf{\ja{ぺらぺら}}, \emph{perapera}.)
\end{tcolorbox}
